\title{GLM: General Language Model Pretraining with Autoregressive Blank Infilling}

\begin{abstract}
There have been various types of pretraining architectures including autoencoding models (e.g., BERT), autoregressive models (e.g., GPT), and encoder-decoder models (e.g., T5). However, none of the pretraining frameworks performs the best for all tasks of three main categories including natural language understanding (NLU), unconditional generation, and conditional generation. We propose a General Language Model (GLM)  based on autoregressive blank infilling to address this challenge. GLM improves blank filling pretraining by adding 2D positional encodings and allowing an arbitrary order to predict spans, which results in performance gains over BERT and T5 on NLU tasks. 

Meanwhile, GLM can be pretrained for different types of tasks by varying the number and lengths of blanks. On a wide range of tasks across NLU, conditional and unconditional generation, GLM outperforms BERT, T5, and GPT given the same model sizes and data, and achieves the best performance from a single pretrained model with 1.25$\times$ parameters of BERT$_\text{Large}$, demonstrating its generalizability to different downstream tasks.\footnote{The code and pre-trained models are available at \url{https://github.com/THUDM/GLM}}
\end{abstract}

\section{Introduction}
\label{sec:intro}
Language models pretrained on unlabeled texts have substantially advanced the state of the art in various NLP tasks, ranging from natural language understanding (NLU) to text generation~\cite{GPT2018a,devlinBERT2019,YangXLNet2019,GPT2-2018,raffelT52020,lewisBART2019,GPT3-2020}. Downstream task performance as well as the scale of the parameters have also constantly increased in the past few years.

In general, existing pretraining frameworks can be categorized into three families: \textit{autoregressive}, \textit{autoencoding}, and \textit{encoder-decoder} models. Autoregressive models, such as GPT~\cite{GPT2018a}, learn left-to-right language models. While they succeed in long-text generation and show few-shot learning ability when scaled to billions of parameters~\cite{GPT2-2018,GPT3-2020}, the inherent disadvantage is the unidirectional attention mechanism, which cannot fully capture the dependencies between the context words in NLU tasks. Autoencoding models, such as BERT~\cite{devlinBERT2019}, learn bidirectional context encoders via denoising objectives, e.g. Masked Language Model (MLM). The encoders produce contextualized representations that suit natural language understanding tasks, but could not be directly applied for text generation. Encoder-decoder models adopt bidirectional attention for the encoder, unidirectional attention for the decoder, and cross attention between them~\cite{songMASS2019,PALM2020,lewisBART2019}. They are typically deployed in conditional generation tasks, such as text summarization and response generation.
\footnote{Unconditional generation refers to generating text as a language model without finetuning, while conditional generation refers to sequence-to-sequence tasks.}. T5~\cite{raffelT52020} unifies NLU and conditional generation via encoder-decoder models but requires more parameters to match the performance of BRET-based models such as RoBERTa~\cite{RoBERTa} and DeBERTa \cite{He2021DeBERTaDB}. 

None of these pretraining frameworks is flexible enough to perform competitively across all NLP tasks. Previous works have tried to unify different frameworks by combining their objectives via multi-task learning~\cite{UniLM19,UniLMv220}. However, since the autoencoding and autoregressive objectives differ by nature, a simple unification cannot fully inherit the advantages of both frameworks.

In this paper, we propose a pretraining framework named GLM (General Language Model), based on autoregressive blank infilling. We randomly blank out continuous spans of tokens from the input text, following the idea of autoencoding, and train the model to sequentially reconstruct the spans, following the idea of autoregressive pretraining (see \ref{fig:example}). While blanking filling has been used in T5~\cite{raffelT52020} for text-to-text pretraining, we propose two improvements, namely span shuffling and 2D positional encoding. Empirically, we show that with the same amount of parameters and computational cost, GLM significantly outperforms BERT on the SuperGLUE benchmark by a large margin of 4.6\% -- 5.0\% and outperforms RoBERTa and BART when pretrained on a corpus of similar size (158GB). GLM also significantly outperforms T5 on NLU and generation tasks with fewer parameters and data.

Inspired by Pattern-Exploiting Training (PET) \cite{PET:2001-07676}, we reformulate NLU tasks as manually-crafted cloze questions that mimic human language. Different from the BERT-based models used by PET, GLM can naturally handle multi-token answers to the cloze question via autoregressive blank filling.

Furthermore, we show that by varying the number and lengths of missing spans, the autoregressive blank filling objective can pretrain language models for conditional and unconditional generation. Through multi-task learning of different pretraining objectives, a single GLM can excel in both NLU and (conditional and unconditional) text generation. Empirically, compared with standalone baselines, GLM with multi-task pretraining achieves improvements in NLU, conditional text generation, and language modeling tasks altogether by sharing the parameters.

\section{GLM Pretraining Framework} We propose a general pretraining framework GLM based on a novel autoregressive blank infilling objective. GLM formulates NLU tasks as cloze questions that contain task descriptions, which can be answered by autoregressive generation. 

\subsection{Pretraining Objective}
\label{subsec:objective}

\subsubsection{Autoregressive Blank Infilling}
\label{subsec:blankinfill}

GLM is trained by optimizing an \textit{autoregressive blank infilling} objective. Given an input text $\boldsymbol{x}=[x_1,\cdots,x_n]$, multiple text spans $\{\boldsymbol{s}_1,\cdots,\boldsymbol{s}_m\}$ are sampled, where each span $\boldsymbol{s}_i$ corresponds to a series of consecutive tokens $[s_{i,1},\cdots,s_{i,l_i}]$ in $\boldsymbol{x}$.  Each span is replaced with a single $[\text{MASK}]$ token, forming a corrupted text $\boldsymbol{x}_{\text{corrupt}}$. The model predicts the missing tokens in the spans from the corrupted text in an autoregressive manner, which means when predicting the missing tokens in a span, the model has access to the corrupted text \emph{and} the previously predicted spans. To fully capture the interdependencies between different spans, we randomly permute the order of the spans, similar to the permutation language model~\cite{YangXLNet2019}. Formally, let $Z_m$ be the set of all possible permutations of the length-$m$ index sequence $[1,2,\cdots,m]$, and $\boldsymbol{s}_{\boldsymbol{z}_{<i}}$ be $[\boldsymbol{s}_{z_1},\cdots,\boldsymbol{s}_{z_{i-1}}]$, we define the pretraining objective as

\begin{equation}
    \max_\theta \mathbb{E}_{\boldsymbol{z}\sim Z_m}\left[\sum_{i=1}^m\log p_\theta(\boldsymbol{s}_{z_i}|\boldsymbol{x}_{\text{corrupt}},\boldsymbol{s}_{\boldsymbol{z}_{<i}})\right]
    \label{eqn:objective}
\end{equation}

We always generate the tokens in each blank following a left-to-right order, i.e. the probability of generating the span $\boldsymbol{s}_i$ is factorized as:

\begin{align}
\begin{split}
    &p_\theta(\boldsymbol{s}_i|\boldsymbol{x}_{\text{corrupt}},\boldsymbol{s}_{\boldsymbol{z}_{<i}})\\
    =&\prod_{j=1}^{l_i} p(s_{i,j}|\boldsymbol{x}_{\text{corrupt}},\boldsymbol{s}_{\boldsymbol{z}_{<i}},\boldsymbol{s}_{i,<j})
\end{split}
\end{align}

We implement the autoregressive blank infilling objective with the following techniques. The input $\boldsymbol{x}$ is divided into two parts: Part A is the corrupted text $\boldsymbol{x}_{\text{corrupt}}$, and Part B consists of the masked spans. Part A tokens can attend to each other, but cannot attend to any tokens in B. Part B tokens can attend to Part A and antecedents in B, but cannot attend to any subsequent tokens in B. To enable autoregressive generation, each span is padded with special tokens $[\text{START}]$ and $[\text{END}]$, for input and output respectively. In this way, our model automatically learns a bidirectional encoder (for Part A) and a unidirectional decoder (for Part B) in a unified model. The implementation of GLM is illustrated in \ref{fig:model}.

We randomly sample spans of length drawn from a Poisson distribution with $\lambda=3$. We repeatedly sample new spans until at least 15\% of the original tokens are masked. Empirically, we have found that the 15\% ratio is critical for good performance on downstream NLU tasks.

\subsubsection{Multi-Task Pretraining}
\label{subsec:multitask}

In the previous section, GLM masks short spans and is suited for NLU tasks. However, we are interested in pretraining a single model that can handle both NLU and text generation. We then study a \textit{multi-task pretraining} setup, in which a second objective of generating longer text is jointly optimized with the blank infilling objective. We consider the following two objectives:

\begin{itemize}
    \item Document-level. We sample a single span whose length is sampled from a uniform distribution over 50\%--100\% of the original length. The objective aims for long text generation.

    \item Sentence-level. We restrict that the masked spans must be full sentences. Multiple spans (sentences) are sampled to cover 15\% of the original tokens. This objective aims for seq2seq tasks whose predictions are often complete sentences or paragraphs. 
\end{itemize}

Both new objectives are defined in the same way as the original objective, i.e. Eq.~\ref{eqn:objective}. The only difference is the number of spans and the span lengths.

\subsection{Model Architecture} GLM uses a single Transformer with several modifications to the architecture: (1) we rearrange the order of layer normalization and the residual connection, which has been shown critical for large-scale language models to avoid numerical errors~\cite{Megatron-LM}; (2) we use a single linear layer for the output token prediction; (3) we replace ReLU activation functions with GeLUs~\cite{GeLU:HendrycksG16}.

\subsubsection{2D Positional Encoding} One of the challenges of the autoregressive blank infilling task is how to encode the positional information. Transformers rely on positional encodings to inject the absolute and relative positions of the tokens. We propose 2D positional encodings to address the challenge. Specifically, each token is encoded with two positional ids. The first positional id represents the position in the corrupted text $\boldsymbol{x}_{\text{corrupt}}$. For the masked spans, it is the position of the corresponding $[\text{MASK}]$ token. The second positional id represents the intra-span position. For tokens in Part A, their second positional ids are $0$. For tokens in Part B, they range from 1 to the length of the span. The two positional ids are projected into two vectors via learnable embedding tables, which are both added to the input token embeddings.

Our encoding method ensures that the model is not aware of the length of the masked span when reconstructing them. It is an important difference as compared to other models. For example, XLNet~\cite{YangXLNet2019} encodes the original position so that it can perceive the number of missing tokens, and SpanBERT~\cite{joshiSpanBERT2020} replaces the span with multiple [MASK] tokens and keeps the length unchanged. Our design fits downstream tasks as usually the length of the generated text is unknown beforehand.

\subsection{Finetuning GLM}
\label{subsec:finetune}
Typically, for downstream NLU tasks, a linear classifier takes the representations of sequences or tokens produced by pretrained models as input and predicts the correct labels.

The practices are different from the generative pretraining task, leading to inconsistency between pretraining and finetuning. 

Instead, we reformulate NLU classification tasks as generation tasks of blank infilling, following PET~\cite{PET:2001-07676}. Specifically, given a labeled example $(\boldsymbol{x}, y)$, we convert the input text $\boldsymbol{x}$ to a cloze question $c(\boldsymbol{x})$ via a pattern containing a single mask token. The pattern is written in natural language to represent the semantics of the task. For example, a sentiment classification task can be formulated as ``\{SENTENCE\}. It's really $[\text{MASK}]$''. The candidate labels $y\in\mathcal{Y}$ are also mapped to answers to the cloze, called verbalizer $v(y)$. In sentiment classification, the labels ``positive'' and ``negative'' are mapped to the words ``good'' and ``bad''. The conditional probability of predicting $y$ given $\boldsymbol{x}$ is

\begin{equation}
    p(y|\boldsymbol{x})=\frac{p(v(y)|c(\boldsymbol{x}))}{\sum_{y'\in \mathcal{Y}}p(v(y')|c(\boldsymbol{x}))}
\end{equation}

where $\mathcal{Y}$ is the label set. Therefore the probability of the sentence being positive or negative is proportional to predicting ``good'' or ``bad'' in the blank. Then we finetune GLM with a cross-entropy loss (see \ref{fig:finetune}).

For text generation tasks, the given context constitutes the Part A of the input, with a mask token appended at the end. The model generates the text of Part B autoregressively. We can directly apply the pretrained GLM for unconditional generation, or finetune it on downstream conditional generation tasks.

\subsection{Discussion and Analysis} In this section, we discuss the differences between GLM and other pretraining models. We are mainly concerned with how they can be adapted to downstream blank infilling tasks.

\textbf{Comparison with BERT~\cite{devlinBERT2019}.}

As pointed out by \cite{YangXLNet2019}, BERT fails to capture the interdependencies of masked tokens due to the independence assumption of MLM. Another disadvantage of BERT is that it cannot fill in the blanks of multiple tokens properly. To infer the probability of an answer of length $l$,  BERT needs to perform $l$ consecutive predictions. If the length $l$ is unknown, we may need to enumerate all possible lengths, since BERT needs to change the number of $[\text{MASK}]$ tokens according to the length.

\textbf{Comparison with XLNet~\cite{YangXLNet2019}.} Both GLM and XLNet are pretrained with autoregressive objectives, but there are two differences between them. First, XLNet uses the original position encodings before corruption. During inference, we need to either know or enumerate the length of the answer, the same problem as BERT. Second, XLNet uses a two-stream self-attention mechanism, instead of the right-shift, to avoid the information leak within Transformer. It doubles the time cost of pretraining.

\textbf{Comparison with T5~\cite{raffelT52020}.} T5 proposes a similar blank infilling objective to pretrain an encoder-decoder Transformer. 

T5 uses independent positional encodings for the encoder and decoder, and relies on multiple sentinel tokens to differentiate the masked spans. In downstream tasks, only one of the sentinel tokens is used, leading to a waste of model capacity and inconsistency between pretraining and finetuning. Moreover, T5 always predicts spans in a fixed left-to-right order. As a result, GLM can significantly outperform T5 on NLU and seq2seq tasks with fewer parameters and data, as stated in \ref{subsec:single_obj,subsec:multi_obj}.

\textbf{Comparison with UniLM~\cite{UniLM19}.} UniLM combines different pretraining objectives under the autoencoding framework by changing the attention mask among bidirectional, unidirectional, and cross attention. However, UniLM always replaces masked spans with [MASK] tokens, which limits its ability to model the dependencies between the masked spans and their context. GLM feeds in the previous token and autoregressively generates the next token. Finetuning UniLM on downstream generation tasks also relies on masked language modeling, which is less efficient. UniLMv2~\cite{UniLMv220} adopts partially autoregressive modeling for generation tasks, along with the autoencoding objective for NLU tasks. Instead, GLM unifies NLU and generation tasks with autoregressive pretraining.

\section{Experiments}
\label{sec:experiment}

We now describe our pretraining setup and the evaluation of downstream tasks.

\subsection{Pretraining Setup}
\label{subsec:data}
For a fair comparison with BERT~\cite{devlinBERT2019}, we use BooksCorpus~\cite{BookCorpusZhu5} and English Wikipedia as our pretraining data. We use the uncased wordpiece tokenizer of BERT with 30k vocabulary. 

We train GLM$_\text{Base}$ and GLM$_\text{Large}$ with the same architectures as BERT$_\text{Base}$ and BERT$_\text{Large}$, containing 110M and 340M parameters respectively. 

For multi-task pretraining, we train two Large-sized models with a mixture of the blank infilling objective and the document-level or sentence-level objective, denoted as GLM$_\text{Doc}$ and GLM$_\text{Sent}$.  Additionally, we train two larger GLM models of 410M (30 layers, hidden size 1024, and 16 attention heads) and 515M (30 layers, hidden size 1152, and 18 attention heads) parameters with document-level multi-task pretraining, denoted as GLM$_\text{410M}$ and GLM$_\text{515M}$.

To compare with SOTA models, we also train a Large-sized model with the same data, tokenization, and hyperparameters as RoBERTa~\cite{RoBERTa}, denoted as GLM$_{\text{RoBERTa}}$.

Due to resource limitations, we only pretrain the model for 250,000 steps, which are half of RoBERTa and BART's training steps and close to T5 in the number of trained tokens. More experiment details can be found in \ref{sec:appendix_pretrain}.

\subsection{SuperGLUE}
\label{subsec:single_obj}

To evaluate our pretrained GLM models, we conduct experiments on the SuperGLUE benchmark~\cite{SuperGLUE2019} and report the standard metrics. SuperGLUE consists of 8 challenging NLU tasks. We reformulate the classification tasks as blank infilling with human-crafted cloze questions, following PET~\cite{PET2:2009-07118}. Then we finetune the pretrained GLM models on each task as described in \ref{subsec:finetune}. The cloze questions and other details can be found in \ref{subsec:appendix_superglue}.

For a fair comparison with GLM$_\text{Base}$ and GLM$_\text{Large}$, we choose BERT$_\text{Base}$ and BERT$_\text{Large}$ as our baselines, which are pretrained on the same corpus and for a similar amount of time. We report the performance of standard finetuning (i.e. classification on the [CLS] token representation). The performance of BERT with cloze questions is reported in \ref{subsec:ablation}. To compare with GLM$_{\text{RoBERTa}}$, we choose T5, BART$_\text{Large}$, and RoBERTa$_\text{Large}$ as our baselines. T5 has no direct match in the number of parameters for BERT$_\text{Large}$, so we present the results of both T5$_\text{Base}$ (220M parameters) and T5$_\text{Large}$ (770M parameters). All the other baselines are of similar size to BERT$_\text{Large}$.

\ref{tab:superglue} shows the results. With the same amount of training data, GLM consistently outperforms BERT on most tasks with either base or large architecture. The only exception is WiC (word sense disambiguation). On average, GLM$_\text{Base}$ scores 4.6\% higher than BERT$_\text{Base}$, and GLM$_\text{Large}$ scores 5.0\% higher than BERT$_\text{Large}$. It clearly demonstrates the advantage of our method in NLU tasks. In the setting of RoBERTa$_\text{Large}$, GLM$_{\text{RoBERTa}}$ can still achieve improvements over the baselines, but with a smaller margin. Specifically, GLM$_{\text{RoBERTa}}$ outperforms T5$_\text{Large}$ but is only half its size. We also find that BART does not perform well on the challenging SuperGLUE benchmark. We conjecture this can be attributed to the low parameter efficiency of the encoder-decoder architecture and the denoising sequence-to-sequence objective. 

\subsection{Multi-Task Pretraining}
\label{subsec:multi_obj}
Then we evaluate the GLM's performance in a multi-task setting (\ref{subsec:objective}). Within one training batch, we sample short spans and longer spans (document-level or sentence-level) with equal chances. We evaluate the multi-task model for NLU, seq2seq, blank infilling, and zero-shot language modeling.

\textbf{SuperGLUE.}
For NLU tasks, we evaluate models on the SuperGLUE benchmark. The results are also shown in \ref{tab:superglue}. We observe that with multi-task pretraining, GLM$_\text{Doc}$ and GLM$_\text{Sent}$ perform slightly worse than GLM$_\text{Large}$, but still outperform BERT$_\text{Large}$ and UniLM$_\text{Large}$. Among multi-task models, GLM$_\text{Sent}$ outperforms GLM$_\text{Doc}$ by 1.1\% on average. Increasing GLM$_\text{Doc}$'s parameters to 410M (1.25$\times$BERT$_\text{Large}$) leads to better performance than GLM$_\text{Large}$. GLM with 515M parameters (1.5$\times$BERT$_\text{Large}$) can perform even better.

\textbf{Sequence-to-Sequence.}
\label{sec:seq2seq}
Considering the available baseline results, we use the Gigaword dataset~\cite{GigawordRushCW15} for abstractive summarization and the SQuAD 1.1 dataset~\cite{SQuAD:RajpurkarZLL16} for question generation~\cite{NQG:du2017} as the benchmarks for models pretrained on BookCorpus and Wikipedia. Additionally, we use the CNN/DailyMail~\cite{GetPointSummarization2017} and XSum~\cite{XSum2018} datasets for abstractive summarization as the benchmarks for models pretrained on larger corpora. 

The results for models trained on BookCorpus and Wikipedia are shown in \ref{tab:gigaword,tab:squad_qg}. We observe that GLM$_\text{Large}$ can achieve performance matching the other pretraining models on the two generation tasks.
GLM$_\text{Sent}$ can perform better than GLM$_\text{Large}$, while
GLM$_\text{Doc}$ performs slightly worse than GLM$_\text{Large}$. This indicates that the document-level objective, which teaches the model to extend the given contexts, is less helpful to conditional generation, which aims to extract useful information from the context. Increasing GLM$_\text{Doc}$'s parameters to 410M leads to the best performance on both tasks. The results for models trained on larger corpora are shown in \ref{tab:cnndm_xsum}. GLM$_{\text{RoBERTa}}$ can achieve performance matching the seq2seq BART model, and outperform T5 and UniLMv2.

\textbf{Text Infilling.} Text infilling is the task of predicting missing spans of text which are consistent with the surrounding context \cite{TextInfill2019,Fillin2020,BLM2020}. GLM is trained with an autoregressive blank infilling objective, thus can straightforwardly solve this task. We evaluate GLM on the Yahoo Answers dataset~\cite{YahooAnswers} and compare it with Blank Language Model (BLM) \cite{BLM2020}, which is a specifically designed model for text infilling. From the results in \ref{tab:yahoo}, GLM outperforms previous methods by large margins (1.3 to 3.9 BLEU) and achieves the state-of-the-art result on this dataset. We notice that GLM$_\text{Doc}$ slightly underperforms GLM$_\text{Large}$, which is consistent with our observations in the seq2seq experiments.

\textbf{Language Modeling.}

Most language modeling datasets such as WikiText103 are constructed from Wikipedia documents, which our pretraining dataset already contains.

Therefore, we evaluate the language modeling perplexity on a held-out test set of our pretraining dataset, which contains about 20M tokens, denoted as BookWiki. We also evaluate GLM on the LAMBADA dataset~\cite{LAMBADADataset2016}, which tests the ability of systems to model long-range dependencies in text. The task is to predict the final word of a passage. As the baseline, we train a GPT$_\text{Large}$ model~\cite{GPT2-2018,GPT3-2020} with the same data and tokenization as GLM$_\text{Large}$. 

The results are shown in \ref{tab:language_model}. All the models are evaluated in the zero-shot setting. Since GLM learns the bidirectional attention, we also evaluate GLM under the setting in which the contexts are encoded with bidirectional attention. Without generative objective during pretraining, GLM$_\text{Large}$ cannot complete the language modeling tasks, with perplexity larger than 100. With the same amount of parameters, GLM$_\text{Doc}$ performs worse than GPT$_\text{Large}$. This is expected since GLM$_\text{Doc}$ also optimizes the blank infilling objective. Increasing the model's parameters to 410M (1.25$\times$ of GPT$_\text{Large}$) leads to a performance close to GPT$_\text{Large}$. GLM$_\text{515M}$ (1.5$\times$ of GPT$_\text{Large}$) can further outperform GPT$_\text{Large}$. With the same amount of parameters, encoding the context with bidirectional attention can improve the performance of language modeling. Under this setting, GLM$_\text{410M}$ outperforms GPT$_\text{Large}$. This is the advantage of GLM over unidirectional GPT. We also study the contribution of 2D positional encoding to long text generation. We find that removing the 2D positional encoding leads to lower accuracy and higher perplexity in language modeling.

\textbf{Summary.} 
Above all, we conclude that GLM effectively shares model parameters across natural language understanding and generation tasks, achieving better performance than a standalone BERT, encoder-decoder, or GPT model. 

\subsection{Ablation Study}
\label{subsec:ablation}

\ref{tab:classifier} shows our ablation analysis for GLM. First, to provide an apple-to-apple comparison with BERT, we train a BERT$_\text{Large}$ model  with our implementation, data, and hyperparameters (row 2). The performance is slightly worse than the official BERT$_\text{Large}$ and significantly worse than GLM$_\text{Large}$. It confirms the superiority of GLM over Masked LM pretraining on NLU tasks. Second, we show the SuperGLUE performance of GLM finetuned as sequence classifiers (row 5) and BERT with cloze-style finetuning (row 3). Compared to BERT with cloze-style finetuning, GLM benefits from the autoregressive pretraining. Especially on ReCoRD and WSC, where the verbalizer consists of multiple tokens, GLM consistently outperforms BERT. This demonstrates GLM's advantage in handling variable-length blank. Another observation is that the cloze formulation is critical for GLM's performance on NLU tasks. For the large model, cloze-style finetuning can improve the performance by 7 points. Finally, we compare GLM variants with different pretraining designs to understand their importance. Row 6 shows that removing the span shuffling (always predicting the masked spans from left to right) leads to a severe performance drop on SuperGLUE. Row 7 uses different sentinel tokens instead of a single $[\text{MASK}]$ token to represent different masked spans. The model performs worse than the standard GLM. We hypothesize that it wastes some modeling capacity to learn the different sentinel tokens which are not used in downstream tasks with only one blank. In \ref{tab:language_model}, we show that removing the second dimension of 2D positional encoding hurts the performance of long text generation.

We note that T5 is pretrained with a similar blank infilling objective. GLM differs in three aspects: (1) GLM consists of a single encoder, (2) GLM shuffles the masked spans, and (3) GLM uses a single [MASK] instead of multiple sentinel tokens. While we cannot directly compare GLM with T5 due to the differences in training data and the number of parameters, the results in \ref{tab:superglue,tab:classifier} have demonstrated the advantage of GLM.

\section{Related Work}
\textbf{Pretrained Language Models.} Pretraining large-scale language models significantly improves the performance of downstream tasks. There are three types of pretrained models. First, autoencoding models learn a bidirectional contextualized encoder for natural language understanding via denoising objectives~\cite{devlinBERT2019,joshiSpanBERT2020,YangXLNet2019,RoBERTa,ALBERT:lan2020,ELECTRA:clark2020}. 

Second, autoregressive models are trained with a left-to-right language modeling objective~\cite{GPT2018a,GPT2-2018,GPT3-2020}. 

Third, encoder-decoder models are pretrained for sequence-to-sequence tasks~\cite{songMASS2019,lewisBART2019,PALM2020,PEGASUS:zhang2020}. 

Among encoder-decoder models, BART~\cite{lewisBART2019} conducts NLU tasks by feeding the same input into the encoder and decoder, and taking the final hidden states of the decoder. Instead, T5~\cite{raffelT52020} formulates most language tasks in the text-to-text framework. However, both models require more parameters to outperform autoencoding models such as RoBERTa~\cite{RoBERTa}. UniLM~\cite{UniLM19,UniLMv220} unifies three pretraining models under the masked language modeling objective with different attention masks.

\textbf{NLU as Generation.} Previously, pretrained language models complete classification tasks for NLU with linear classifiers on the learned representations. GPT-2~\cite{GPT2-2018} and GPT-3~\cite{GPT3-2020} show that generative language models can complete NLU tasks such as question answering by directly predicting the correct answers without finetuning, given task instructions or a few labeled examples. However, generative models require much more parameters to work due to the limit of unidirectional attention. Recently, PET~\cite{PET:2001-07676,PET2:2009-07118} proposes to reformulate input examples as cloze questions with patterns similar to the pretraining corpus in the few-shot setting. It has been shown that combined with gradient-based finetuning, PET can achieve better performance in the few-shot setting than GPT-3 while requiring only 0.1\% of its parameters. Similarly, \citet{Augmented2020} and \citet{Augmented2021} convert structured prediction tasks, such as sequence tagging and relation extraction, to sequence generation tasks. 

\textbf{Blank Language Modeling.} \citet{Fillin2020} and \citet{BLM2020} also study blanking infilling models. Different from their work, we pre-train language models with blank infilling objectives and evaluate their performance in downstream NLU and generation tasks.

\section{Conclusions}
GLM is a general pretraining framework for natural language understanding and generation. We show that the NLU tasks can be formulated as conditional generation tasks, and therefore solvable by autoregressive models. GLM unifies the pretraining objectives for different tasks as autoregressive blank infilling, with mixed attention masks and the novel 2D position encodings. Empirically we show that GLM outperforms previous methods for NLU tasks and can effectively share parameters for different tasks. 

\end{document}
